\documentclass[12pt]{article}

\usepackage[spanish]{babel}
\usepackage[utf8]{inputenc}
\usepackage[T1]{fontenc}
\usepackage{mathpazo}
\usepackage{amsmath, amssymb, amsfonts}
\usepackage{mathtools}
\usepackage{geometry}
\usepackage{xcolor}
\usepackage{graphicx}
\usepackage{tikz}
\usetikzlibrary{babel, arrows.meta}
\usepackage{enumitem}
\usepackage{multicol}
\usepackage{array}
\usepackage{booktabs}
\usepackage[skins,breakable]{tcolorbox}
\usepackage{fancyhdr}
\usepackage{titlesec}
\usepackage{hyperref}

\geometry{letterpaper, margin=2.3cm}
\setlength{\parindent}{0pt}
\setlength{\parskip}{6pt}
\linespread{1.06}

\definecolor{tinta}{RGB}{38, 38, 42}
\definecolor{acento}{RGB}{110, 45, 50}
\definecolor{acentosuave}{RGB}{249, 244, 238}
\definecolor{gris}{RGB}{105, 105, 105}
\definecolor{grissuave}{RGB}{247, 246, 244}
\definecolor{linea}{RGB}{195, 190, 180}

\hypersetup{colorlinks=true, linkcolor=tinta, urlcolor=acento, citecolor=acento}

\setlength{\headheight}{14pt}
\pagestyle{fancy}
\fancyhf{}
\lhead{\small\scshape\color{gris} Clase de Física --- Joaquín Molina}
\rhead{\small\color{gris} \thepage}
\renewcommand{\headrulewidth}{0.3pt}
\renewcommand{\headrule}{\color{linea}\hrule height\headrulewidth}

\titleformat{\section}
    {\Large\scshape\color{tinta}}
    {\thesection}
    {0.6em}
    {}
    [{\color{linea}\vspace{-0.4em}\hrule height 0.4pt}]
\titleformat{\subsection}
    {\large\itshape\color{tinta}}
    {\thesubsection}
    {0.5em}
    {}

\tcbset{
    enhanced, breakable, frame hidden, sharp corners,
    left=3mm, right=3mm, top=1.5mm, bottom=1.5mm, boxsep=0.5mm,
    fonttitle=\scshape
}
\newtcolorbox{definicion}[1][]{colback=white,
    borderline west={1.6pt}{0pt}{acento}, coltitle=acento,
    title=Definición, #1}
\newtcolorbox{ejemplo}[1][]{colback=grissuave,
    borderline west={1.6pt}{0pt}{gris}, coltitle=gris,
    title=Ejemplo, #1}
\newtcolorbox{observacion}[1][]{colback=white,
    borderline west={1pt}{0pt}{linea}, coltitle=gris,
    fonttitle=\itshape,
    title=Observación, #1}
\newtcolorbox{importante}[1][]{colback=acentosuave,
    borderline west={2.2pt}{0pt}{tinta}, coltitle=tinta,
    fonttitle=\scshape\bfseries,
    title=Importante, #1}
\newtcolorbox{ejercicio}[1][]{colback=white,
    borderline west={1.6pt}{0pt}{acento}, coltitle=acento,
    title=Ejercicio, #1}

\newcommand{\R}{\mathbb{R}}
\newcommand{\N}{\mathbb{N}}
\newcommand{\Z}{\mathbb{Z}}
\newcommand{\Q}{\mathbb{Q}}
\newcommand{\abs}[1]{\left|#1\right|}
\newcommand{\paren}[1]{\left(#1\right)}
\newcommand{\corchete}[1]{\left[#1\right]}

\newcommand{\tituloClase}{Fundamento Teórico detrás del VQE}
\newcommand{\subtituloClase}{Taller de bienvenida --- Grupo MoleQLs}
\newcommand{\nombreProfesor}{Joaquín Molina}

\begin{document}

\begin{center}
    {\Huge\scshape \tituloClase}\\[0.35cm]
    {\color{linea}\rule{0.35\textwidth}{0.4pt}}\\[0.35cm]
    {\large\itshape \subtituloClase}\\[0.25cm]
    {\normalsize \nombreProfesor}
\end{center}

\vspace{0.6cm}
{\color{linea}\hrule height 0.4pt}
\vspace{0.6cm}

\section{Antes de empezar: qué asumiremos conocido}

Esta clase busca llegar directo al algoritmo VQE (\textit{Variational Quantum
Eigensolver}), así que no partiremos desde cero absoluto. Para que el tiempo
rinda, vamos a asumir que ya manejas los siguientes conceptos:

\begin{importante}[title=\textbf{Prerrequisitos asumidos}]
\begin{itemize}
    \item Qué es un \textbf{espacio de Hilbert} y por qué los estados
    cuánticos viven ahí.
    \item La notación de \textbf{ket} $|\psi\rangle$ y \textbf{bra}
    $\langle\psi|$, y el producto interno $\langle\phi|\psi\rangle$.
    \item Qué es un \textbf{qubit}: $|\psi\rangle = \alpha|0\rangle +
    \beta|1\rangle$ con $|\alpha|^2 + |\beta|^2 = 1$.
    \item Sistemas de \textbf{varios qubits} vía producto tensorial
    ($|q_1\rangle \otimes |q_2\rangle \otimes \dots$).
    \item \textbf{Compuertas cuánticas básicas}: $X$, $H$ (Hadamard),
    $R_y(\theta)$, $R_z(\theta)$ y CNOT.
\end{itemize}
\end{importante}

Si algo de esta lista no te suena familiar, no te preocupes: es normal, y
mejor repasarlo antes de seguir para que el resto de la clase tenga sentido.

\begin{observacion}[title=\textbf{Libros recomendados para estos prerrequisitos}]
\begin{itemize}
    \item \textbf{M. Nielsen \& I. Chuang}, \textit{Quantum Computation and
    Quantum Information}. El libro de referencia del área completa; los
    capítulos 1 y 2 cubren exactamente todo lo asumido arriba.
    \item \textbf{J.J. Sakurai}, \textit{Modern Quantum Mechanics}. Ideal si
    quieres entender bien el formalismo de espacios de Hilbert, kets/bras y
    operadores desde el lado de la mecánica cuántica ``tradicional''.
    \item \textbf{R. Sutor}, \textit{Dancing with Qubits}. Buena introducción
    suave, con harta intuición antes de la matemática formal.
\end{itemize}
\end{observacion}

\newpage

\section{El Hamiltoniano como operador de energía}

\begin{definicion}
El \textbf{Hamiltoniano} $\hat{H}$ de un sistema cuántico es el operador
Hermítico ($\hat{H} = \hat{H}^\dagger$) asociado a su energía total. Al ser
Hermítico, sus autovalores son siempre reales, y corresponden a los
posibles valores de energía que se pueden medir en el sistema.
\end{definicion}

La ecuación de autoestados que usaremos en todo lo que sigue no es un
postulado aparte: se obtiene directamente de la ecuación de Schrödinger
dependiente del tiempo.

\begin{ejemplo}[title=\textbf{Demostración: de Schrödinger a la ecuación de autoestados}]
Partimos de la ecuación de Schrödinger dependiente del tiempo:
\[
i\hbar \frac{\partial}{\partial t}|\Psi(t)\rangle = \hat{H} |\Psi(t)\rangle
\]

\begin{enumerate}
    \item Si $\hat{H}$ no depende explícitamente del tiempo, buscamos
    soluciones donde la parte espacial y la temporal se separan:
    $|\Psi(t)\rangle = |\phi\rangle \, f(t)$.

    \item Sustituimos en la ecuación de Schrödinger:
    \[
    i\hbar \, |\phi\rangle \, \frac{df}{dt} = f(t) \, \hat{H} |\phi\rangle
    \]

    \item Dividimos ambos lados por el \textbf{escalar} $f(t)$ (esto sí es
    válido: $f(t) \in \mathbb{C}$, no es un vector, así que dividir por él
    tiene sentido; nunca dividimos por el ket $|\phi\rangle$):
    \[
    \hat{H} |\phi\rangle = \left[\frac{i\hbar}{f(t)}\frac{df}{dt}\right]
    |\phi\rangle
    \]

    \item El lado izquierdo, $\hat{H}|\phi\rangle$, no depende de $t$
    (ni $\hat{H}$ ni $|\phi\rangle$ dependen del tiempo). Para que la
    igualdad se cumpla en \textit{todo} instante, el escalar entre
    corchetes del lado derecho está obligado a ser constante en el tiempo.
    Llamamos a esa constante $E$:
    \[
    \frac{i\hbar}{f(t)}\frac{df}{dt} = E
    \qquad \text{y} \qquad
    \hat{H} |\phi\rangle = E |\phi\rangle
    \]

    \item La primera ecuación da $f(t) = e^{-iEt/\hbar}$: solo una fase
    global, así que $|\Psi(t)\rangle = |\phi\rangle e^{-iEt/\hbar}$ no
    cambia sus probabilidades en el tiempo (\textit{estado estacionario}).

    \item La segunda ecuación, $\hat{H}|\phi\rangle = E|\phi\rangle$, es
    exactamente la \textbf{ecuación de autoestados}: la constante de
    separación $E$ resulta ser un autovalor de $\hat{H}$, y por eso se
    interpreta como una energía del sistema. $\blacksquare$
\end{enumerate}
\end{ejemplo}

En general $\hat{H}$ tiene varios autoestados, que indexamos con $n$:

\[
\hat{H} |\phi_n\rangle = E_n |\phi_n\rangle
\]

Cada $|\phi_n\rangle$ es un \textbf{autoestado} de $\hat{H}$, y $E_n$ su
\textbf{autovalor} asociado: la energía que se mide si el sistema se
encuentra en ese estado.

\begin{importante}
Al conjunto ordenado de energías $E_0 \leq E_1 \leq E_2 \leq \dots$ se le
llama \textbf{espectro} de $\hat{H}$. El menor autovalor $E_0$ es la
\textbf{energía del estado fundamental} (\textit{ground state}), y su
autoestado asociado $|\phi_0\rangle$ es el estado de mínima energía del
sistema.
\end{importante}

\begin{ejemplo}[title=\textbf{Ejemplo: espectro de un Hamiltoniano de un qubit}]
Sea $\hat{H} = Z$ (Pauli-Z). Como $Z|0\rangle = |0\rangle$ y $Z|1\rangle =
-|1\rangle$, los estados $|0\rangle$ y $|1\rangle$ son autoestados de
$\hat{H}$ con autovalores $E=+1$ y $E=-1$ respectivamente. Como $-1 < 1$,
el estado fundamental de este Hamiltoniano es $|\phi_0\rangle = |1\rangle$,
con energía $E_0 = -1$.
\end{ejemplo}

En este sentido, ``resolver'' un sistema cuántico equivale a diagonalizar
$\hat{H}$: encontrar todos sus autovalores y autoestados. En química
cuántica, en particular, encontrar $E_0$ es casi siempre lo que más
interesa, pues determina propiedades como la estabilidad de una molécula o
su geometría de equilibrio.

Diagonalizar $\hat{H}$ exactamente, sin embargo, puede ser muy costoso.
Nos gustaría contar con una forma de \textbf{aproximar} $E_0$ sin tener que
diagonalizar. La siguiente herramienta matemática es justo eso, y es el
corazón de VQE.

\newpage

\section{El principio variacional}

\begin{importante}[title=\textbf{Principio variacional}]
Para cualquier estado normalizado $|\psi\rangle$ y cualquier Hamiltoniano
$\hat{H}$ con estado fundamental de energía $E_0$:
\[
\langle \psi | \hat{H} | \psi \rangle \geq E_0
\]
con igualdad si y solo si $|\psi\rangle$ es el estado fundamental.
\end{importante}

En otras palabras: \textbf{cualquier estado que probemos nos da una cota
superior de $E_0$}. Si vamos ajustando el estado para minimizar
$\langle\hat{H}\rangle$, nos vamos acercando a la energía real del estado
fundamental.

\subsection{De autovalores a un problema de minimización}

Buscar $E_0$ diagonalizando $\hat{H}$ significa revisar, en principio,
\textbf{todo} el espacio de Hilbert. El principio variacional nos permite
cambiar de estrategia: en vez de eso, restringimos la búsqueda a una
familia de estados de prueba, \textbf{parametrizada} por un conjunto de
números reales $\vec\theta = (\theta_1, \theta_2, \dots)$, que denotamos
$|\psi(\vec\theta)\rangle$. Con esto definimos una función de energía que
depende solo de los parámetros:
\[
E(\vec\theta) = \langle \psi(\vec\theta) | \hat{H} | \psi(\vec\theta) \rangle
\]

\begin{importante}
Por el principio variacional, $E(\vec\theta) \geq E_0$ para \textbf{todo}
$\vec\theta$. Por lo tanto:
\[
E_0 \approx \min_{\vec\theta} \; E(\vec\theta)
\]
Encontrar el estado fundamental (un problema de autovalores) se transforma
en \textbf{minimizar una función real de varias variables} (un problema de
optimización clásica): mientras menor sea $E(\vec\theta)$, más cerca está
$|\psi(\vec\theta)\rangle$ de ser el autoestado fundamental $|\phi_0\rangle$.
\end{importante}

Esta es la idea completa detrás de VQE: convertir el problema cuántico de
diagonalizar $\hat{H}$ en un problema de optimización sobre $\vec\theta$,
donde un computador clásico puede encargarse de la minimización. Falta
responder algo concreto: ¿cómo se construye $|\psi(\vec\theta)\rangle$ en
la práctica?

\newpage

\section{Cómo parametrizar un estado}

Un computador cuántico solo puede transformar estados mediante
\textbf{compuertas unitarias}: operadores $U$ tales que $U^\dagger U = I$
(preservan la norma del estado, como debe ser). Para construir
$|\psi(\vec\theta)\rangle$ necesitamos compuertas que además dependan de un
parámetro real $\theta$ que podamos ajustar continuamente.

\begin{definicion}
La forma estándar de construir una \textbf{compuerta parametrizada} es como
la exponencial de un operador $\hat{G}$:
\[
U(\theta) = e^{-i\theta \hat{G}}
\]
A $\hat{G}$ se le llama el \textbf{generador} de la compuerta.
\end{definicion}

No cualquier operador $\hat{G}$ sirve: para que $U(\theta)$ sea unitaria
para todo $\theta \in \R$, el generador $\hat{G}$ tiene que ser
\textbf{Hermítico}.

\begin{ejemplo}[title=\textbf{Demostración: por qué $\hat{G}$ debe ser Hermítico}]
Queremos que $U(\theta) = e^{-i\theta\hat{G}}$ cumpla $U^\dagger(\theta)
\, U(\theta) = I$.

\begin{enumerate}
    \item Tomamos el adjunto de $U(\theta)$. El adjunto de una exponencial
    de operador conjuga tanto el escalar como el operador:
    \[
    U^\dagger(\theta) = \left(e^{-i\theta\hat{G}}\right)^{\!\dagger}
    = e^{\,i\theta \hat{G}^\dagger}
    \]

    \item Entonces:
    \[
    U^\dagger(\theta)\, U(\theta) = e^{\,i\theta \hat{G}^\dagger} \,
    e^{-i\theta \hat{G}}
    \]

    \item Si $\hat{G}$ es Hermítico ($\hat{G}^\dagger = \hat{G}$), los dos
    exponentes tienen el \textbf{mismo operador} (y por lo tanto conmutan
    entre sí), así que podemos sumarlos directamente:
    \[
    e^{\,i\theta \hat{G}} \, e^{-i\theta \hat{G}}
    = e^{\,i\theta\hat{G} - i\theta\hat{G}} = e^{0} = I
    \]

    \item Si $\hat{G}$ \textbf{no} fuera Hermítico, $\hat{G}^\dagger \neq
    \hat{G}$, los exponentes no se cancelarían, y $U(\theta)$ no
    preservaría la norma del estado. $\blacksquare$
\end{enumerate}
\end{ejemplo}

\begin{importante}
$U(\theta) = e^{-i\theta\hat{G}}$ es unitaria para todo $\theta$
\textbf{si y solo si} $\hat{G}$ es Hermítico. Esta es la razón por la que
las compuertas cuánticas parametrizadas siempre se construyen a partir de
operadores Hermíticos: los mismos operadores que usamos como observables
(sección 2) sirven como generadores de compuertas.
\end{importante}

\begin{ejemplo}[title=\textbf{Ejemplo: rotaciones de un qubit}]
Los operadores de Pauli $X$, $Y$, $Z$ son Hermíticos, así que generan
compuertas unitarias válidas. Las rotaciones de un qubit que ya conocías
son exactamente de esta forma:
\[
R_x(\theta) = e^{-i\theta X/2}, \qquad
R_y(\theta) = e^{-i\theta Y/2}, \qquad
R_z(\theta) = e^{-i\theta Z/2}
\]
\end{ejemplo}

Con compuertas de este tipo (y CNOTs para entrelazar qubits), el estado
parametrizado completo se construye como una cadena de exponenciales
aplicadas al estado inicial $|0\rangle^{\otimes n}$:
\[
|\psi(\vec\theta)\rangle = \left(\prod_k e^{-i\theta_k \hat{G}_k}\right)
|0\rangle^{\otimes n}
\]
donde cada $\hat{G}_k$ es Hermítico (un Pauli o un Pauli string), y cada
$\theta_k$ es un parámetro ajustable que el optimizador clásico de VQE va
modificando.

\end{document}
